% Source: https://artofproblemsolving.com/community/c3h2502367_basic_problems_on_function
% Stable identity: aops-topic-2502367-q01
\textbf{One solution.} Let $g(x)=\dfrac{(x-1)x(x+1)(x+2)}{4}$. Then
\[g(x+1)-g(x)=\dfrac{x(x+1)(x+2)\big[(x+3)-(x-1)\big]}{4}=x(x+1)(x+2).\]

\textbf{All solutions.} If $f$ is any polynomial solution, then $h=f-g$ is a polynomial with $h(x+1)=h(x)$ for all $x$. In particular $h(n)=h(0)$ for every positive integer $n$, so the polynomial $h(x)-h(0)$ has infinitely many roots and is the zero polynomial. Hence $h$ is constant, and
\[\boxed{f(x)=\dfrac{(x-1)x(x+1)(x+2)}{4}+C}\]
for a constant $C$; all of these work.
